\section{What has been established, and what has not}
\label{sec:status}

The sections above reported results of very different standing: identities verified in exact
arithmetic, numerical results with a stated accuracy, consequences of additional assumptions, and
estimates. This section sorts them. An item is called established if a machine check stands behind
it (Appendix~\ref{app:checks}); statements of interpretation deserve the reader's independent
scrutiny more than the numbers do.

\subsection{Established}

\emph{Geometry} (Sec.~\ref{sec:geometry}). The curvature and the Einstein tensor of the deflating
spacetime in closed form, for arbitrary rates and constants. The Killing vectors, 14, 15 or 20 in
number, and the completeness of the lists. That the spacetime is not a solution of Einstein's
equations without matter, with or without a cosmological constant; that with three equal rates it is
an exact solution of Einstein-Gauss-Bonnet gravity without matter, for $\Lam=3\,\gff^{2}$ and
$\aGB=-1/(4\,\gff^{2})$; that with unequal rates it is not; and that this solution is unstable
(Sec.~\ref{sec:gravity}).
\SIG{The turning value of the ordinary time for a classical particle with momentum along a deflating
time.}{The unbounded blueshift of a classical particle that moves along a deflating coordinate, and
the incompleteness of its geodesic.}

\emph{The Dirac field} (Secs.~\ref{sec:dirac}, \ref{sec:modes}). The Dirac matrices in coordinate
components, the spin connection, the covariant constancy of the matrices, and the curvature of the
spinor derivative. The Lagrangian with the kinetic term, the mass term and the self-interaction; the
field equation (\ref{eq:diracexp}) with the self-interaction on its right-hand side; the weighted
field; and the statement that fields which do not depend on $x5,x6,x7$ obey the flat equation with
the couplings $g_r/\Vd$ of the moment.
\SIG{The exact solutions when one rate governs; that exactly eight of the sixteen solutions are
bounded; that these are standing waves with vanishing charge; and the law of their extinction,
Eqs.~(\ref{eq:magnitude}) to (\ref{eq:constant}).}{The exact solutions when one rate governs, and
the law of creation (\ref{eq:ncr}) with the temperature (\ref{eq:Tdefl}).} All of this is confirmed
numerically by two independent methods.

\emph{The self-interaction} (Sec.~\ref{sec:selfint}). The Fierz matrix in closed form and from the
256 matrices. That for anticommuting fields the nine invariants obey five linear relations, so that
only four couplings are independent, verified by exact integer arithmetic and on Grassmann-valued
fields. The four effective couplings (\ref{eq:geff}), the Hartree-Fock couplings (\ref{eq:GHF}), the
pair couplings (\ref{eq:Gpair}), and the conditions under which the self-interaction keeps the chiral
symmetry and the symmetry $U(1)_E$.\SIG{ The exact relation (\ref{eq:forms}) between the two forms of
the self-interaction in the quantized sector.}{}

\emph{Matter with the self-interaction} (Sec.~\ref{sec:ks}). The Kohn-Sham states of uniform matter
and their thermodynamic consistency.\SIG{}{ The reference calculation on flat spacetime, reproduced
to ten digits by two programs. The history of matter with momentum along the deflating coordinates:
the growth of its density, the fall of the mass of its particles, the limited time during which it
is self-bound, the anisotropy of its pressure, and the conservation law (\ref{eq:conservation}).}
The thresholds of Table~\ref{tab:ksthr} at which a gas becomes self-bound matter in the sector of
fields that do not depend on the deflating coordinates, and the history of the saturated matter
after that moment.\SIG{ That in the covariant form with the positive kinetic term the two halves of
the field polarize.}{}

\emph{The effective potential with the self-interaction} (Sec.~\ref{sec:cjtcontact}). That the
Kohn-Sham method and the CJT effective potential at the order of the Hartree and Fock diagrams give
the same states. The critical coupling (\ref{eq:Gc}) of the vacuum with the filled states of
negative energy, and the moment at which the deflation crosses it. The functional
(\ref{eq:unified}) with bare mass, self-interaction and Vorton-Scalar; that the exchange of the
radial mode acts at low density as a scalar contact coupling; the shifts of the moment of binding of
Table~\ref{tab:unified}; and the mass (\ref{eq:mtheta}) that a bare Dirac mass gives to the phase
mode.

\emph{Pairing with the self-interaction} (Sec.~\ref{sec:gapcontact}). The amplitudes
(\ref{eq:Fclosed}) of the four contact channels in closed form; that the exchange of the
Vorton-Scalar is a self-interaction of finite range; and the gaps of Table~\ref{tab:gapcontact}.

\emph{Gravity} (Sec.~\ref{sec:gravity}). The Einstein-Lovelock tensors (\ref{eq:E1}) and
(\ref{eq:E2}), derived in three ways. The reference evolution, reproduced by two programs. That the
undisturbed deflation needs the constants (\ref{eq:tuned}) or a source with the equation of state
(\ref{eq:agent}); that the matter of the model cannot be that source; that matter accelerates the
deflation and makes the flat coordinates expand; and that the undisturbed deflation is unstable in
the tuned gravity.

\emph{The model along the clock} (Secs.~\ref{sec:clock} to \ref{sec:cjt}). For fields that do not
depend on the extra coordinates, the model on the deflating spacetime is the model on flat spacetime
with the couplings of the moment, exactly. The masses are constant; $\gVS$ grows as
$e^{\td/2}$; $\aVS$, $\GVS$, $\DVS$ and $C$ grow as $e^{\td}$. The modes of the Vorton-Scalar grow
with the weight only, and the potential between two fermions has the unchanged range and the
strength of the moment. The exchange amplitudes are those of flat spacetime times the square of the
coupling of the moment. In weak coupling the gap grows according to Eq.~(\ref{eq:growth}), and the
critical temperature follows it. Two fermions become bound at a definite time, at a scattering
resonance, with a binding energy that starts quadratically; the order of the other events depends on
$m_s/m$ alone. The paired state follows the deflation if and only if $\Xid$ is below about 0.2 to
0.5; before that it lags and accumulates broken pairs. The solution of the gap equation is a minimum
of the effective potential at every moment examined. Matter of fixed density loses its fermion mass
at a critical value of $C$, is self-bound for $C>3.7011$ and is more strongly bound later; the
melting temperature of the fermion mass falls; the barrier of the one-loop potential of the vacuum
moves toward the vacuum value.

\emph{Order parameter and symmetry} (Sec.~\ref{sec:order}). The potential of the order parameter
prefers no direction among the four frame axes. The internal symmetry of the sector, the
representation of the pair operators and the breaking pattern are unchanged by the deflation. Of the
six generators that belong to the pairs of the four extra axes, the deflating spacetime keeps none,
one or three as symmetries of the spacetime. A vortex line of unit winding carries four fermion zero
modes, two for each direction along the line, at every moment; its energy per length falls as
$e^{-\td}$.

\emph{Cosmology} (Sec.~\ref{sec:cosmo}), each item under the assumption named there. The equation of
the expansion (\ref{eq:friedmann}) and the fall of the inferred Hubble constant, with recombination
computed on the deflating background and verified by an independent calculation. The bounds of
Table~\ref{tab:bounds}, and that in the model as it stands the atomic clocks give none. The effect
of the loss of mass of bound pairs. The pressure (\ref{eq:pagent}) of the agent. The equations
(\ref{eq:pert}) of the perturbations of ordinary space for unperturbed extra scale factors, the
identities (\ref{eq:id1}) and (\ref{eq:id2}), and that the four equations cannot all hold with matter
alone. The growth equations of the two closures, and that the closure ``Shear'' coincides with the
closure ``Pressure'' inside the horizon. The spectra of the microwave background for $\ell\ge30$ and
their independence of the closure; the forecast for the rate from mock data; that the Hubble
constant falls in every fit, with or without extra radiation; and that extra radiation and the rate
of deflation compensate in part in the spectra, without any gain for the Hubble constant. The
unchanged estimates of degeneracy
and of condensation in halos, the limits on superfluid dark matter, the epoch of the resonance, and
the drift of the quantities.\SIG{}{ The temperature of the created quanta at the allowed rates.}

\emph{The Higgs hypothesis} (Sec.~\ref{sec:higgs}). Under H2 the electroweak scale follows the
coupling with the exponent $\nEW$, and atomic clocks bound the rate of deflation by
Eq.~(\ref{eq:clockbound}). The exponent is given in closed form for the scale-invariant ladder
equation.

\subsection{What the mass term and the self-interaction change}
\label{sec:changes}

Sections~\ref{sec:gap} to \ref{sec:order} were derived for the model in which the whole mass of the
fermions comes from the Vorton-Scalar and the whole interaction from its exchange ($m_D=0$,
$g_r=0$). Table~\ref{tab:changes} lists, result by result, what a bare Dirac mass and the
self-interaction change.

\begin{table*}
\caption{\label{tab:changes}What a bare Dirac mass $m_D$ and the self-interaction (couplings $g_r$)
change in the results obtained with the Vorton-Scalar alone.}
\begin{ruledtabular}
\begin{tabular}{p{0.27\textwidth}p{0.33\textwidth}p{0.33\textwidth}}
result with the Vorton-Scalar alone & with a bare Dirac mass & with the self-interaction \\
\hline
fields that do not depend on $x5,x6,x7$ obey the flat equations with couplings of the moment (Secs.~\ref{sec:dirac}, \ref{sec:clock}) & unchanged; $m_D$ itself does not change with the clock & unchanged; every $g_r$ grows as $e^{\td}$, in a fixed ratio to $\gVS^{2}$ \\
the chiral symmetry $\UVS$ is exact and the phase mode massless (Sec.~\ref{sec:clock}) & broken explicitly; the phase mode has the mass of Eq.~(\ref{eq:mtheta}), which grows as $e^{\td/2}$ & kept only if $\geff0+\geff8=0$ \\
the gap of the paired state and its growth (Sec.~\ref{sec:gap}) & smaller, because the massive phase mode binds less & larger with a scalar contact attraction, smaller with a vector repulsion (Table~\ref{tab:gapcontact}) \\
two fermions bind when $\DVS$ reaches $D_c$ (Sec.~\ref{sec:pair}) & the Yukawa potential of the radial mode is unchanged & a contact term $\Gpair$ is added to the potential; the threshold with it has not been computed \\
matter is self-bound for $C>3\pi^{2}/8$ (Sec.~\ref{sec:cjt}) & earlier (Table~\ref{tab:unified}) & later with a vector repulsion; earlier with a scalar attraction where the particles keep a mass (Table~\ref{tab:unified}) \\
matter of fixed density loses its whole fermion mass (Sec.~\ref{sec:cjt}) & the particles keep the mass $m_D$ & the mass in matter is $\mst$ of Eq.~(\ref{eq:hKS}) \\
the internal symmetry $[SU(2)\times SU(2)]\times U(1)_E$ of the sector (Sec.~\ref{sec:order}) & unchanged & $U(1)_E$ is broken by vector and axial-vector couplings \\
cosmology and the bounds on the rate (Sec.~\ref{sec:cosmo}) & unchanged & unchanged \\
\end{tabular}
\end{ruledtabular}
\end{table*}

\subsection{The two methods as checks of each other and of the other results}
\label{sec:crosschecks}

The Kohn-Sham method and the CJT effective potential were used not only for new results but as
checks, and each result of this paper that could be computed in two ways was
(Table~\ref{tab:crosschecks}).

\begin{table*}
\caption{\label{tab:crosschecks}Results obtained in two or more independent ways, with the largest
relative deviation between them.}
\begin{ruledtabular}
\begin{tabular}{p{0.34\textwidth}p{0.46\textwidth}c}
result & obtained by & deviation \\
\hline
{{ROWS:tbl_crosschecks}}
\end{tabular}
\end{ruledtabular}
\end{table*}

\subsection{Not established}

\emph{The source of the deflation.} The metric is prescribed. For equal rates it solves the
equations of Einstein-Gauss-Bonnet gravity without matter, but that solution is unstable
(Sec.~\ref{sec:gravity}); in Einstein gravity it needs a source with the equation of state
(\ref{eq:agent}), which the matter of the model cannot be. For unequal rates no field equation has
been found that the metric solves. What holds the deflation, and how it absorbs the reaction of
matter, is not known. All results on fields are results on a given background.

\SIG{\emph{The form of the self-interaction on the $4+4$ spacetime.} The covariant Lagrangian has an
indefinite kinetic term in the quantized sector, and the self-interaction can be completed in two
ways (Sec.~\ref{sec:forms}). They differ in the sign of the interaction between particles of
different halves of the field, and they give different matter: the quantized-sector form binds it,
the covariant form polarizes it and binds it much less. This paper adopts the quantized-sector form
because the model with the Vorton-Scalar is defined in it; no principle has been found that selects
it.

}{}\emph{The cutoff.} The self-interaction has couplings of negative mass dimension and defines an
effective theory. The results without the filled states of negative energy presuppose a cutoff below
the mass of the particles (Sec.~\ref{sec:sea}); with the filled states the results depend on the
cutoff, and the mass of the phase mode depends on it through the scalar density of the vacuum. The
cutoff is a parameter of the model that nothing in the hypotheses fixes.

\emph{The couplings.} The nine couplings $g_r$, of which four are independent, are free parameters.
The examples use a scalar attraction and a vector repulsion in the ratio 10 to 1; no value is
predicted.\SIG{}{ For the reference example the two numbers $-10^{5}$ and $+10^{4}$ are taken as the
mean-field couplings themselves; taken as couplings of the Lagrangian they would not bind matter
(Sec.~\ref{sec:ksflat}).}

\emph{Matter beyond the mean field.} The Kohn-Sham states and the effective potential at the order
of the Hartree and Fock diagrams neglect correlations. Without the filled states the energy has no
lowest value as a function of an unrestricted trial mass; the states reported are solutions of the
gap equation. Temperature was not included in the calculations with the self-interaction.\SIG{}{ The
matter with momentum along the deflating coordinates was treated without collisions; whether
collisions are frequent enough to release the heat of Sec.~\ref{sec:ks71} was not examined.}

\emph{Pairing inside dense matter.} The gaps of Table~\ref{tab:gapcontact} are computed for
particles of the mass $m$. In self-bound matter the mass is $\mst<m$; the gap equation and the
equation for $\mst$ have not been solved together. The contact couplings enter the gap equation with
a range $1/\mct$ that is a choice.

\emph{The theory of gravity and the assumption A2.} If gravity is of the Einstein-Gauss-Bonnet kind,
the relation (\ref{eq:GN}) between Newton's constant and the volume of the deflating coordinates,
and with it every result of Sec.~\ref{sec:cosmo} that is marked A2, would have to be derived anew.

\emph{The agent.} Nothing in the hypothesis says what keeps the three rates constant while ordinary
space expands. Under A2 it must have the tension (\ref{eq:pagent}) and no energy density; whether a
physical source with these properties exists has not been examined. Nor does the hypothesis say how
the agent responds when ordinary space is perturbed. That the scale factors of the deflating
coordinates are not perturbed (assumption A3) is an assumption; and even with it, a closure is
needed that the hypothesis does not supply. The growth of structure, the amplitude $\sigma_8$ and
the spectra of the microwave background at the largest angular scales are computed for two closures
and differ between them at the first order in the rate. They are the quantities that a theory of
the agent would have to predict.

%%<1644>
\emph{The exclusion of the growing solutions.} The eight growing solutions of each mode that depends
on a deflating time are discarded because they are unbounded. That is a condition of admissibility;
it is not derived from a principle. The theory has no positive conserved charge for such modes, and
a consistent quantization exists only for the sector that does not depend on the additional times.
%%</1644>
%%<1671>
\emph{The reaction of the created particles on the metric.} The particles created by the deflation
gain energy without bound. Their stress must eventually change the metric; this has not been
computed. Nor do the coordinates cover the whole future of particles that move along deflating
coordinates; what lies beyond is outside the hypothesis.

\emph{The restriction to the sector.} That the fields do not depend on $x5,x6,x7$ is assumed
wherever the sector is used. The deflation supports the assumption only if the deflating coordinates
are periodic (Sec.~\ref{sec:modes}).
%%</1671>

\emph{Unequal rates with simultaneous crossings.} No closed form is known for a mode that depends on
two deflating coordinates whose crossings coincide; the results are numerical.

\emph{The hidden coordinate.} That the fields do not depend on $x0$ is assumed. Nothing in the
deflation explains it.

\emph{The end of the history.} The couplings grow without bound. The four-dimensional description
with the exchange of one boson holds only while $\aVS$ is not large, and the eight-dimensional theory
itself is an effective theory, with couplings of negative mass dimension, that must fail when the
proper size of the deflating coordinates reaches its scale of validity. What happens then is outside
the model. Every history of Secs.~\ref{sec:gap} to \ref{sec:cjt} ends in a regime that the two-loop
truncation does not control; the numbers at the last moments of each table indicate the trend and
are not predictions.

\emph{The energy.} The energy of matter is not conserved while the couplings change. Where it goes,
into whatever drives the deflation, is not described.

\emph{Dynamics.} Sections~\ref{sec:gap}, \ref{sec:pair} and \ref{sec:cjt} give sequences of
equilibrium states. Only for the paired state has the ability to follow been examined, in a model
without dissipation; with dissipation the broken pairs would relax, and their lifetime is not known.
The time-dependent two-body problem at the moment of binding, which decides how much of the
population ends in the bound state, has not been solved. How a gas reaches the self-bound state, and
whether the abnormal state forms by nucleation, are open. The equilibrium of a loop of vortex line
with its current, and its stability, have not been computed.

\emph{The vacuum beyond one loop.} The statements about the vacuum use the one-loop potential
outside its range of reliability and draw only the conclusion stated in Sec.~\ref{sec:cjt}.

\emph{Cosmology.} The forecasts of Sec.~\ref{sec:forecast} are for mock data without gravitational
lensing; no measured spectrum of the microwave background has been fitted. All bounds that reach
back to nucleosynthesis or to recombination assume a constant rate; a rate that was large early and
is small today is neither excluded nor described. Only linear perturbations were computed. The
consequences of the epoch of the resonance for the formation of halos have not been simulated.

\emph{The Higgs hypothesis.} That a
condensate forms in the vacuum, that \SIG{any symmetry of the model}{the rotation group of the extra
dimensions} is the electroweak gauge group, the scale and the mass of the scalar particle, and the
masses of ordinary fermions: none of these is established. The exponent $\nEW$ has not been computed for the model itself; the
statement that it is of order one or larger rests on the scale-invariant case and on the general
fact that the scale of the condensate vanishes at the critical coupling.

\section{Conclusions}
\label{sec:concl}

We asked what becomes of a 16-component Dirac field with a mass term and a self-interaction, and of
the superfluid of its pairs, when three of the eight coordinates of a spacetime of signature
\SIG{$4+4$}{$7+1$} deflate exponentially.

The geometry is simple and rigid. The deflating spacetime is homogeneous, it has 14 to 20 symmetries
instead of 36, and it is not a solution of Einstein's equations without matter. With three equal
rates it is an exact solution of Einstein-Gauss-Bonnet gravity without matter, for one relation
between the rate and the two constants of that theory; but the solution is unstable, the matter of
the model cannot be the source that Einstein gravity would need, and matter that is present
accelerates the deflation.

The Dirac field separates into two parts. Fields that do not depend on the deflating coordinates
obey, after a change of normalization, exactly the flat equations.
\SIG{Fields that do depend on them are extinguished: every admissible mode is a standing wave that
dies out as the exponential of an exponential of the time. The sector in which the $4+4$ theory is
not unitary is removed by the deflation, although the condition that selects the admissible modes is
imposed and not derived.}{Fields that do depend on them are populated by the deflation with a
Fermi-Dirac spectrum at the temperature of the rate divided by $2\pi$, and their quanta then gain
energy without bound.}

The self-interaction has nine channels and, because the field anticommutes, four independent
couplings. On the deflating spacetime they grow with the inverse volume of the deflating
coordinates, in step with the coupling of the Vorton-Scalar. The Kohn-Sham method and the effective
action of Cornwall, Jackiw and Tomboulis give the same matter, and the deflation carries it through
definite events: a gas becomes self-bound, a massless field acquires a mass, and with a bare Dirac
mass the phase mode of the Vorton-Scalar becomes heavy.\SIG{ On the $4+4$ spacetime the
self-interaction must be defined in the quantized sector, and its two possible forms give different
matter.}{ Matter with momentum along the deflating coordinates is compressed, loses the mass of its
particles, and is self-bound only for about one e-folding of the deflating volume.}

The superfluid of hypothesis H1 exists on the deflating spacetime at every moment, with the
couplings of that moment, and it acquires a history. The masses do not change; the coupling grows
with the inverse volume of the deflating coordinates; the gap grows much faster than the coupling,
in weak coupling as the exponential of an exponential; two fermions become bound at a scattering
resonance; matter loses its fermion mass and
becomes self-bound; and the paired state can follow the deflation only once its gap exceeds a few
times the rate divided by the pairing strength. A scalar contact attraction raises the gap, and a
vector repulsion lowers it. The direction of the order parameter is not selected
by the deflation.

For hypotheses H2 and H3 the deflation brings no relief. \SIG{The Higgs hypothesis had no geometric
basis on the flat $4+4$ spacetime and has none here.}{The rotations that gave the Higgs hypothesis
its one structural support on the flat $7+1$ spacetime are no longer symmetries of the spacetime.}
If the superfluid were the Higgs field, the electroweak scale would follow the growing coupling, and
atomic clocks would then limit the rate of deflation to about a millionth of the Hubble rate. If
gravity is the Einstein theory in eight dimensions, lunar laser ranging limits it to three
thousandths of the Hubble rate, and the Hubble constant inferred from the microwave background
moves away from the local measurement, in the simplest estimate and in fits of all parameters to
mock data alike. In either case nothing of the history just described can have happened within the
age of the universe at a constant rate.

The perturbations of ordinary space show most plainly what the hypothesis of deflation leaves open.
Whatever keeps the rates constant must carry a tension, and it must take part in the perturbations:
with a Newton's constant that depends on the time, the Einstein equations of ordinary space cannot
all hold with matter alone. The spectra of the microwave background at all but the largest angular
scales do not depend on how the agent does so, and mock data of the quality of the Planck satellite
would bound the rate to four per cent of the Hubble rate. The growth of structure does depend on it,
at the first order in the rate.

These results are quantitative, and they come with the programs that produced them. They mark what
a successful theory of this kind would have to supply: a source or a theory of gravity that sustains
the deflation stably and that says how the source responds to matter and to perturbations, a rate that was large early
and is small today, and a reason why the constants of ordinary matter do not follow the volume of
the deflating coordinates.

\section*{Data availability}

The source code of all programs, the data files from which every figure and every computed number of
this paper is generated, the instructions for compiling and running the programs
(Appendix~\ref{app:checks}), and the records of the provenance of each result are openly available
\cite{DeflationRepo}.

